\documentclass{article}
\usepackage{amsmath}
\usepackage{braket}
\usepackage{geometry}
\geometry{letterpaper, margin=1in}

\title{Mathematical Derivation of the Perturbed IQFT Circuit}
\author{}
\date{}

\begin{document}

\maketitle

\section{The Input Encoding}
Let $n$ be the total number of qubits ($n=4$ for this implementation). A classical integer input $x \in \{0, \dots, 2^n-1\}$ corresponds to a multi-qubit computational basis state. Defining $x_k$ as the $k$-th bit of the binary representation of $x$, the initial quantum state is encoded as:
\begin{equation}
    \ket{\psi_0} = \ket{x} = \ket{x_{n-1} x_{n-2} \dots x_1 x_0}
\end{equation}

\section{The Inverse QFT Representation}
The Quantum Fourier Transform operations naturally factorize into independent single-qubit superpositions. When applying the Inverse QFT (IQFT) to an input state $\ket{x}$, it generates a uniform superposition over all possible classical states $\ket{y}$ but rotates their local phases using negative imaginary exponentials. 

According to Qiskit's ``little-endian'' convention (where $q_0$ serves as the least significant bit), the mathematically factorized tensor product of the localized IQFT output is:
\begin{equation}
    \ket{\psi_{\text{iqft}}} = \text{IQFT}\ket{x} = \bigotimes_{k=0}^{n-1} \frac{1}{\sqrt{2}} \left( \ket{0}_k + \exp\left(-2\pi i \frac{x}{2^{n-k}}\right) \ket{1}_k \right)
\end{equation}

\section{Applying The Parameterized Phase Perturbations}
The circuit natively applies parameterized diagonal Phase gates, $P(\theta_k)$, dynamically to every parameterized qubit $k$. A phase gate leaves the $\ket{0}$ amplitude unaffected and rotates the $\ket{1}$ phase component:
\begin{equation}
    P(\theta) \ket{0} = \ket{0}, \quad P(\theta) \ket{1} = e^{i\theta} \ket{1}
\end{equation}

By applying an exclusive phase shift $\theta_k$ to the $k$-th qubit in the factorized sequence, the arbitrary angles geometrically couple to the fractional periods dictated by the IQFT:
\begin{equation}
    \ket{\psi_{\text{out}}} = \bigotimes_{k=0}^{n-1} \frac{1}{\sqrt{2}} \left( \ket{0}_k + \exp\left(-2\pi i \frac{x}{2^{n-k}} + i\theta_k\right) \ket{1}_k \right)
\end{equation}

\section{Global Statevector Equation (In the Computational Basis)}
While the tensor product accurately defines the architecture layer-by-layer, when performing topological measurements or statistical covariance analysis, it is often necessary to isolate the state back into the global $\ket{y}$ superposition form to derive probability densities.

By collapsing the tensor product outward and letting string $y = \sum_{k=0}^{n-1} y_k 2^k$, the local single-qubit phases are aggregated into a single global exponential. The fully expanded mathematical mapping generated by the simulation behaves as follows:
\begin{equation}
    \ket{\psi_{\text{out}}} = \frac{1}{\sqrt{2^n}} \sum_{y=0}^{2^n-1} \exp\left(-2\pi i \frac{xy}{2^n} + i \sum_{k=0}^{n-1} \theta_k y_k \right) \ket{y}
\end{equation}

\section{Appended Inverse Fourier Transform (Final Decoding)}
Applying a secondary Inverse Quantum Fourier Transform to the perturbed phase-encoded state $\ket{\psi_{\text{out}}}$ effectively decodes it, forcing the phase $\theta_k$ parameters to dynamically interfere with the baseline periodicity induced by $x$. 

When applying the final IQFT to every computational state vector $\ket{y}$, each vector transforms into a new superposition over $\ket{z}$. Plugging $\text{IQFT}\ket{y} = \frac{1}{\sqrt{2^n}} \sum_{z=0}^{2^n-1} \exp\left(-2\pi i \frac{yz}{2^n}\right)\ket{z}$ into the global equation yields the final topological state:
\begin{equation}
    \ket{\psi_{\text{final}}} = \text{IQFT} \ket{\psi_{\text{out}}} = \frac{1}{2^n} \sum_{z=0}^{2^n-1} \sum_{y=0}^{2^n-1} \exp\left(-2\pi i \frac{y(x + z)}{2^n} + i \sum_{k=0}^{n-1} \theta_k y_k \right) \ket{z}
\end{equation}

This defines the full complex topological interference pattern. The coefficient magnitude $c_z = \braket{z | \psi_{\text{final}}}$ dictates the structural landscapes observed in the parameter-driven covariance plots.

\section{Marginal Probability of Exact Qubit Measurements}
To physically evaluate the parameterized probability of measuring a singular specific qubit $j$ in the excited $\ket{1}$ state, we isolate its marginal probability density. Because the final IQFT acts globally and deeply entangles the network based on the arbitrary perturbations $\theta_k$, the probability does not cleanly factorize but is instead calculated by tracing out the remaining $(n-1)$ qubits.

Let $z_j \in \{0, 1\}$ denote the binary value of the $j$-th qubit measured in the global collapsed string $z$. The probability $P(q_j = 1)$ is the sum of the squared amplitudes of every basis state $\ket{z}$ where the $j$-th bit is strictly $1$:
\begin{equation}
    P(q_j = 1) = \sum_{z=0}^{2^n-1} z_j \left| c_z \right|^2 
\end{equation}

Substituting the previously derived interference coefficient $c_z$ models the complete marginal expected topography plotted directly on the Z-axis of the 3D surface grid:
\begin{equation}
    P(q_j = 1) = \frac{1}{2^{2n}} \sum_{z=0}^{2^n-1} z_j \left| \sum_{y=0}^{2^n-1} \exp\left(-2\pi i \frac{y(x + z)}{2^n} + i \sum_{k=0}^{n-1} \theta_k y_k \right) \right|^2
\end{equation}

\end{document}
